% Options for packages loaded elsewhere
% Options for packages loaded elsewhere
\PassOptionsToPackage{unicode}{hyperref}
\PassOptionsToPackage{hyphens}{url}
\PassOptionsToPackage{dvipsnames,svgnames,x11names}{xcolor}
%
\documentclass[
  letterpaper,
  DIV=11,
  numbers=noendperiod]{scrreprt}
\usepackage{xcolor}
\usepackage{amsmath,amssymb}
\setcounter{secnumdepth}{5}
\usepackage{iftex}
\ifPDFTeX
  \usepackage[T1]{fontenc}
  \usepackage[utf8]{inputenc}
  \usepackage{textcomp} % provide euro and other symbols
\else % if luatex or xetex
  \usepackage{unicode-math} % this also loads fontspec
  \defaultfontfeatures{Scale=MatchLowercase}
  \defaultfontfeatures[\rmfamily]{Ligatures=TeX,Scale=1}
\fi
\usepackage{lmodern}
\ifPDFTeX\else
  % xetex/luatex font selection
\fi
% Use upquote if available, for straight quotes in verbatim environments
\IfFileExists{upquote.sty}{\usepackage{upquote}}{}
\IfFileExists{microtype.sty}{% use microtype if available
  \usepackage[]{microtype}
  \UseMicrotypeSet[protrusion]{basicmath} % disable protrusion for tt fonts
}{}
\makeatletter
\@ifundefined{KOMAClassName}{% if non-KOMA class
  \IfFileExists{parskip.sty}{%
    \usepackage{parskip}
  }{% else
    \setlength{\parindent}{0pt}
    \setlength{\parskip}{6pt plus 2pt minus 1pt}}
}{% if KOMA class
  \KOMAoptions{parskip=half}}
\makeatother
% Make \paragraph and \subparagraph free-standing
\makeatletter
\ifx\paragraph\undefined\else
  \let\oldparagraph\paragraph
  \renewcommand{\paragraph}{
    \@ifstar
      \xxxParagraphStar
      \xxxParagraphNoStar
  }
  \newcommand{\xxxParagraphStar}[1]{\oldparagraph*{#1}\mbox{}}
  \newcommand{\xxxParagraphNoStar}[1]{\oldparagraph{#1}\mbox{}}
\fi
\ifx\subparagraph\undefined\else
  \let\oldsubparagraph\subparagraph
  \renewcommand{\subparagraph}{
    \@ifstar
      \xxxSubParagraphStar
      \xxxSubParagraphNoStar
  }
  \newcommand{\xxxSubParagraphStar}[1]{\oldsubparagraph*{#1}\mbox{}}
  \newcommand{\xxxSubParagraphNoStar}[1]{\oldsubparagraph{#1}\mbox{}}
\fi
\makeatother


\usepackage{longtable,booktabs,array}
\newcounter{none} % for unnumbered tables
\usepackage{calc} % for calculating minipage widths
% Correct order of tables after \paragraph or \subparagraph
\usepackage{etoolbox}
\makeatletter
\patchcmd\longtable{\par}{\if@noskipsec\mbox{}\fi\par}{}{}
\makeatother
% Allow footnotes in longtable head/foot
\IfFileExists{footnotehyper.sty}{\usepackage{footnotehyper}}{\usepackage{footnote}}
\makesavenoteenv{longtable}
\usepackage{graphicx}
\makeatletter
\newsavebox\pandoc@box
\newcommand*\pandocbounded[1]{% scales image to fit in text height/width
  \sbox\pandoc@box{#1}%
  \Gscale@div\@tempa{\textheight}{\dimexpr\ht\pandoc@box+\dp\pandoc@box\relax}%
  \Gscale@div\@tempb{\linewidth}{\wd\pandoc@box}%
  \ifdim\@tempb\p@<\@tempa\p@\let\@tempa\@tempb\fi% select the smaller of both
  \ifdim\@tempa\p@<\p@\scalebox{\@tempa}{\usebox\pandoc@box}%
  \else\usebox{\pandoc@box}%
  \fi%
}
% Set default figure placement to htbp
\def\fps@figure{htbp}
\makeatother





\setlength{\emergencystretch}{3em} % prevent overfull lines

\providecommand{\tightlist}{%
  \setlength{\itemsep}{0pt}\setlength{\parskip}{0pt}}



 


\KOMAoption{captions}{tableheading}
\makeatletter
\@ifpackageloaded{bookmark}{}{\usepackage{bookmark}}
\makeatother
\makeatletter
\@ifpackageloaded{caption}{}{\usepackage{caption}}
\AtBeginDocument{%
\ifdefined\contentsname
  \renewcommand*\contentsname{Table of contents}
\else
  \newcommand\contentsname{Table of contents}
\fi
\ifdefined\listfigurename
  \renewcommand*\listfigurename{List of Figures}
\else
  \newcommand\listfigurename{List of Figures}
\fi
\ifdefined\listtablename
  \renewcommand*\listtablename{List of Tables}
\else
  \newcommand\listtablename{List of Tables}
\fi
\ifdefined\figurename
  \renewcommand*\figurename{Figure}
\else
  \newcommand\figurename{Figure}
\fi
\ifdefined\tablename
  \renewcommand*\tablename{Table}
\else
  \newcommand\tablename{Table}
\fi
}
\@ifpackageloaded{float}{}{\usepackage{float}}
\floatstyle{ruled}
\@ifundefined{c@chapter}{\newfloat{codelisting}{h}{lop}}{\newfloat{codelisting}{h}{lop}[chapter]}
\floatname{codelisting}{Listing}
\newcommand*\listoflistings{\listof{codelisting}{List of Listings}}
\makeatother
\makeatletter
\makeatother
\makeatletter
\@ifpackageloaded{caption}{}{\usepackage{caption}}
\@ifpackageloaded{subcaption}{}{\usepackage{subcaption}}
\makeatother
\usepackage{bookmark}
\IfFileExists{xurl.sty}{\usepackage{xurl}}{} % add URL line breaks if available
\urlstyle{same}
\hypersetup{
  pdftitle={Rekayasa Cerdas Triune-Intelligence},
  pdfauthor={Norah Jones},
  colorlinks=true,
  linkcolor={blue},
  filecolor={Maroon},
  citecolor={Blue},
  urlcolor={Blue},
  pdfcreator={LaTeX via pandoc}}


\title{Rekayasa Cerdas Triune-Intelligence}
\author{Norah Jones}
\date{2026-04-28}
\begin{document}
\maketitle

\renewcommand*\contentsname{Table of contents}
{
\setcounter{tocdepth}{2}
\tableofcontents
}

\bookmarksetup{startatroot}

\chapter*{Preface}\label{preface}
\addcontentsline{toc}{chapter}{Preface}

\markboth{Preface}{Preface}

This is a Quarto book.

To learn more about Quarto books visit
\url{https://quarto.org/docs/books}.

\bookmarksetup{startatroot}

\chapter{Paradigma, Metodologi, dan Hirarki Rekayasa
TISE}\label{paradigma-metodologi-dan-hirarki-rekayasa-tise}

\section{Pendahuluan: Panggilan Rekayasa Abad
ke-21}\label{pendahuluan-panggilan-rekayasa-abad-ke-21}

Rekayasa pada abad ke-21 menghadapi krisis fundamental berupa devaluasi
agensi manusia akibat otomatisasi yang bersifat menggantikan. Sebagai
respons, paradigma \emph{Triune-Intelligence Smart-Engineering} (TISE)
hadir bukan sebagai evolusi teknis semata, melainkan sebagai panggilan
bagi para rekayasawan untuk menjadi seorang ``Vokator''---problem solver
kreatif yang berdiri teguh di persimpangan antara sains murni dan
filosofi kemanusiaan yang mendalam.

TISE menawarkan antitesis terhadap otomatisasi yang memarginalkan
manusia. Jika rekayasa tradisional sering kali bertujuan menggantikan
peran manusia demi efisiensi, TISE menerapkan prinsip \emph{Leverage \&
Amplification}. Visi ini dapat dipahami melalui analogi ``Penyanyi Rock
dan \emph{Sound System}''. Dalam paradigma lama, mesin dirancang untuk
menyanyi menggantikan sang artis. Sebaliknya, dalam sistem TISE,
teknologi berperan sebagai sistem tata suara raksasa yang tidak
menggantikan suara penyanyi, melainkan mengamplifikasi pesona suara
tersebut agar getarannya mampu menjangkau ratusan ribu pendengar di
stadion global. TISE mengubah teknologi dari alat otomatisasi yang
mematikan peran (displace) menjadi instrumen amplifikasi yang
memberdayakan (empower).

``People Aspire, You Inspire, Engineers Deliver.''

Visi akhir TISE adalah mentransformasi manusia dari konsumen pasif
menjadi agen aktif yang menjalankan misi uniknya. Rekayasa TISE adalah
sebuah janji akademik dan kemanusiaan: untuk tidak membiarkan satu pun
talenta manusia di bumi ini terbuang sia-sia.

\section{Struktur Hirarki Keilmuan
Penopang}\label{struktur-hirarki-keilmuan-penopang}

Penerapan TISE bertumpu pada tiga pilar keilmuan yang saling mengunci
untuk memberikan arah, presisi, dan ``jiwa'' bagi setiap solusi
teknologi:

\textbf{Fundamen (Ilmu Pengetahuan Alam):} Menyediakan hukum-hukum dasar
mengenai materi dan energi sebagai tumpuan realitas fisik dan batas
keterlaksanaan teknis.

\textbf{Penengah (Ilmu Formal/Matematika):} Berfungsi sebagai jembatan
pemroses data yang presisi, logika sistemik, dan pengelola kompleksitas
melalui pemodelan matematis.

\textbf{Penarik/Konteks (Ilmu Kemanusiaan dan Etika):} Sosiologi,
filosofi, dan etika yang memberikan makna, tujuan, dan batasan moral
agar teknologi senantiasa memuliakan martabat manusia sebagai subjek
utama.

\section{Evolusi Rekayasa: Hirarki 7 Level
TISE}\label{evolusi-rekayasa-hirarki-7-level-tise}

Hirarki TISE disusun secara sistematis dari lapisan fisik hingga
ekosistem kemanusiaan global, di mana level yang lebih tinggi memberikan
konteks dan arah, sementara level yang lebih rendah memberikan tumpuan
realitas dan keandalan.

\subsection{Level 0: Lapisan Rekayasa Tradisional (EMI
Core)}\label{level-0-lapisan-rekayasa-tradisional-emi-core}

Pada level ini, fokus rekayasa adalah konversi Energi, Material, dan
Informasi (EMI) melalui \emph{Core Engine}. Berdasarkan
fungsionalitasnya, level ini membedakan antara \textbf{Mesin Statik}
yang berfokus pada stabilitas dan kemampuan menahan beban, serta
\textbf{Mesin Dinamik} yang berfokus pada pengendalian gerak dan aliran.

\textbf{Tabel 3.1: Komponen Core Engine TISE Level 0}

{\def\LTcaptype{none} % do not increment counter
\begin{longtable}[]{@{}
  >{\raggedright\arraybackslash}p{(\linewidth - 2\tabcolsep) * \real{0.5000}}
  >{\raggedright\arraybackslash}p{(\linewidth - 2\tabcolsep) * \real{0.5000}}@{}}
\toprule\noalign{}
\begin{minipage}[b]{\linewidth}\raggedright
Komponen
\end{minipage} & \begin{minipage}[b]{\linewidth}\raggedright
Fungsi Utama
\end{minipage} \\
\midrule\noalign{}
\endhead
\bottomrule\noalign{}
\endlastfoot
Intake & Gerbang kontrol pertama bagi masuknya sumber daya EMI ke dalam
sistem. \\
Encoder & Mengubah sumber daya mentah menjadi bentuk yang siap diproses
(kondisi aliran/sinyal). \\
Decoder & Menerjemahkan hasil pemrosesan internal menjadi aksi atau
keluaran nyata. \\
Exhaust & Mekanisme pelepasan hasil samping atau limbah untuk menjaga
efisiensi siklus. \\
Flywheel & Penyimpan energi kerja yang menjaga kontinuitas, momentum,
dan kestabilan siklus. \\
\end{longtable}
}

\subsection{Level 1: Lapisan Smart Engineering (PSKVE \&
Energon)}\label{level-1-lapisan-smart-engineering-pskve-energon}

Level 1 memperluas domain rekayasa ke arah konversi
\emph{Energon}---segala sesuatu yang memiliki kapasitas untuk melakukan
kerja. Rekayasa ini menggunakan metafora lingkungan kerja yang terdiri
dari \textbf{Stasiun} (titik layanan), \textbf{Jalan} (jalur
distribusi), dan \textbf{Kendaraan} (penggerak/protokol). Mesin
\textbf{PSKVE} bekerja melalui \textbf{Instrumen Konversi} (mengubah
energon menjadi bentuk kerja) dan \textbf{Instrumen Transaksi}
(mengelola aliran nilai).

Lima jenis \emph{Energon} yang dimobilisasi adalah:

\textbf{Energon-1:} Energi fisik dan material tradisional.

\textbf{Energon-2:} Waktu dan perhatian manusia.

\textbf{Energon-3:} Skill, keahlian, dan pengetahuan.

\textbf{Energon-4:} Token bernilai (finansial/aset).

\textbf{Energon-5:} Peran, status, dan keanggotaan sosial.

Mesin ini mengonversi energon menjadi kapasitas kerja dalam lima bentuk:
\emph{Product, Service, Knowledge, Value,} dan \emph{Environment}.

\subsection{Level 2: Lapisan PUDAL Triune (Kecerdasan
Agen)}\label{level-2-lapisan-pudal-triune-kecerdasan-agen}

Level ini mengorkestrasi \emph{Triune Intelligence} (NI, CI, AI) dalam
siklus agentik adaptif menggunakan \emph{Natural Language Prompting}
sebagai medium interaksi lintas kecerdasan.

\textbf{Perception:} NI menangkap intuisi dan emosi; CI menangkap norma
komunitas; AI menangkap data digital dan pola statistik.

\textbf{Understanding:} NI menafsirkan niat; CI memberikan konteks
sosial; AI mengolah klasifikasi dan prediksi data.

\textbf{Decision:} NI memberikan pertimbangan etis; CI memberikan
legitimasi kolektif; AI mensimulasikan konsekuensi logis.

\textbf{Action:} Dilaksanakan melalui aksi manusia (NI), koordinasi
kelompok (CI), atau otomasi sistem (AI).

\textbf{Learning:} NI belajar melalui refleksi; CI melalui audit sosial;
AI melalui penyesuaian model data operasional.

\subsection{Level 3: Lapisan Narasi (Identitas \&
Agensi)}\label{level-3-lapisan-narasi-identitas-agensi}

Sesuai dengan monograf TISE, Level 3 berfokus pada \textbf{Identitas
Naratif} sebagai sumber utama formulasi misi. Manusia dipandang sebagai
protagonis yang menuliskan kisah hidupnya. Rekayasa di sini membangun
``Teater Solusi''---lingkungan di mana stakeholder mementaskan talenta
uniknya. Penggunaan \textbf{Prompt Reflektif} di level ini membantu
stakeholder untuk tidak sekadar menjadi operator teknis, melainkan
memahami peran, nilai, dan panggilan naratifnya dalam sistem.

\subsection{Level 4: Lapisan Meta-TISE DNA (Platform \&
Reproduksi)}\label{level-4-lapisan-meta-tise-dna-platform-reproduksi}

Level 4 adalah rekayasa atas artefak TISE itu sendiri agar menjadi
platform reproduktif. Fokusnya adalah menjaga \textbf{DNA Desain} (pola,
aturan, dan logika operasi) yang memungkinkan lahirnya \textbf{Artefak
Anak}. Proses ini melibatkan ``perkawinan'' antara DNA platform yang
sudah teruji dengan DNA gagasan baru, memastikan prinsip
\emph{reusability} sehingga solusi tidak perlu dibangun dari nol,
melainkan diwariskan dengan adaptasi konteks yang segar.

\subsection{Level 5: Lapisan Valorize (Ekonomi Misi \&
Stakeholder)}\label{level-5-lapisan-valorize-ekonomi-misi-stakeholder}

Level 5 adalah level operasional tertinggi yang memobilisasi stakeholder
bermisi melalui mesin \textbf{Matrix of Mission}. Struktur ini
memastikan adanya pertukaran nilai yang adil dan berkelanjutan.

\textbf{Tabel 3.2: Matrix of Mission (Arsitektur Kontribusi)}

{\def\LTcaptype{none} % do not increment counter
\begin{longtable}[]{@{}
  >{\raggedright\arraybackslash}p{(\linewidth - 8\tabcolsep) * \real{0.2000}}
  >{\raggedright\arraybackslash}p{(\linewidth - 8\tabcolsep) * \real{0.2000}}
  >{\raggedright\arraybackslash}p{(\linewidth - 8\tabcolsep) * \real{0.2000}}
  >{\raggedright\arraybackslash}p{(\linewidth - 8\tabcolsep) * \real{0.2000}}
  >{\raggedright\arraybackslash}p{(\linewidth - 8\tabcolsep) * \real{0.2000}}@{}}
\toprule\noalign{}
\begin{minipage}[b]{\linewidth}\raggedright
Stakeholder
\end{minipage} & \begin{minipage}[b]{\linewidth}\raggedright
USER
\end{minipage} & \begin{minipage}[b]{\linewidth}\raggedright
SOURCE
\end{minipage} & \begin{minipage}[b]{\linewidth}\raggedright
REGULATOR
\end{minipage} & \begin{minipage}[b]{\linewidth}\raggedright
PROVIDER
\end{minipage} \\
\midrule\noalign{}
\endhead
\bottomrule\noalign{}
\endlastfoot
USER & Misi \& Benefit & Kontribusi U ke S & Kontribusi U ke R &
Kontribusi U ke P \\
SOURCE & Kontribusi S ke U & Misi \& Benefit & Kontribusi S ke R &
Kontribusi S ke P \\
REGULATOR & Kontribusi R ke U & Kontribusi R ke S & Misi \& Benefit &
Kontribusi R ke P \\
PROVIDER & Kontribusi P ke U & Kontribusi P ke S & Kontribusi P ke R &
Misi \& Benefit \\
\end{longtable}
}

Keberhasilan level ini ditentukan oleh \textbf{Stakeholder Readiness
Framework} (Arsitektur Kesiapan) yang mencakup 6 dimensi:
\emph{Narrative, Mission, Competency, Motivational, Commitment,} dan
\emph{Opportunity Readiness}.

\subsection{Level 6: Ekosistem Keunggulan (Splendid Theaters of
Life)}\label{level-6-ekosistem-keunggulan-splendid-theaters-of-life}

Puncak evolusi TISE adalah terwujudnya \emph{Splendid Theaters of Life}.
Ini adalah visi akhir di mana 7 miliar manusia mementaskan talenta
uniknya di panggung dunia tanpa degradasi agensi. Dunia bertransformasi
menjadi stadion global digital di mana teknologi memastikan setiap
kontribusi individu memberikan nilai maksimal bagi kemanusiaan secara
keseluruhan.

\section{Metodologi Operasional
Integratif}\label{metodologi-operasional-integratif}

Operasionalisasi TISE mengintegrasikan fase linear rekayasa dengan
siklus operasional \textbf{MOS-7} (\emph{Mission Operating System})
untuk menjaga integritas misi dari awal hingga akhir.

\textbf{Tabel 4.1: Workflow Rekayasa TISE Terintegrasi}

{\def\LTcaptype{none} % do not increment counter
\begin{longtable}[]{@{}
  >{\raggedright\arraybackslash}p{(\linewidth - 8\tabcolsep) * \real{0.2000}}
  >{\raggedright\arraybackslash}p{(\linewidth - 8\tabcolsep) * \real{0.2000}}
  >{\raggedright\arraybackslash}p{(\linewidth - 8\tabcolsep) * \real{0.2000}}
  >{\raggedright\arraybackslash}p{(\linewidth - 8\tabcolsep) * \real{0.2000}}
  >{\raggedright\arraybackslash}p{(\linewidth - 8\tabcolsep) * \real{0.2000}}@{}}
\toprule\noalign{}
\begin{minipage}[b]{\linewidth}\raggedright
Fase Linear
\end{minipage} & \begin{minipage}[b]{\linewidth}\raggedright
Fokus Utama
\end{minipage} & \begin{minipage}[b]{\linewidth}\raggedright
MOS-7 Dominan
\end{minipage} & \begin{minipage}[b]{\linewidth}\raggedright
Level TISE
\end{minipage} & \begin{minipage}[b]{\linewidth}\raggedright
Keluaran Utama
\end{minipage} \\
\midrule\noalign{}
\endhead
\bottomrule\noalign{}
\endlastfoot
Persepsi Masalah & Mendengar keluhan stakeholder & Hear, Model & Level 5
& Rumusan persoalan \& alternatif konfigurasi \\
Konsep \& Model & Memilih konfigurasi terbaik & Hear, Model, Design,
Commit & Level 5, 3 & Matrix of Mission terpilih \& arah narasi \\
Desain \& Planning & Menurunkan kebutuhan teknis & Design, Commit, Build
& Level 4, 2, 1 & DNA Platform \& skema PUDAL/PSKVE \\
Konstruksi & Membangun elemen \& artefak & Build, Commit, Operate &
Level 4 s.d. 0 & Artefak anak \& instrumen layanan hidup \\
Operasional & Menjalankan solusi nyata & Operate, Evaluate & Seluruh
Level & Pengalaman user \& operasi provider nyata \\
Evaluasi & Audit etika \& pembelajaran & Evaluate, Hear, Model & Level 5
& Hasil audit misi \& desain siklus baru \\
\end{longtable}
}

\section{Sintesis dan Kesimpulan}\label{sintesis-dan-kesimpulan}

Hirarki TISE (0-6) membentuk satu rantai rekayasa utuh yang menyatukan
visi kemanusiaan dengan keandalan teknologi. TISE menolak pemisahan
antara efisiensi mesin dan martabat manusia, menempatkan pertumbuhan
manusia sebagai bagian inheren dari proses rekayasa itu sendiri.

\textbf{3 Prinsip Utama Keberhasilan TISE:}

\textbf{Sentralitas Misi:} Solusi harus berakar pada misi naratif
stakeholder, bukan sekadar fungsi teknis.

\textbf{Amplifikasi Agensi:} Teknologi wajib memperkuat dan
mengamplifikasi, bukan menggantikan, talenta unik manusia.

\textbf{Kesiapan Stakeholder (Readiness):} Keberhasilan sistem
bergantung pada pertumbuhan kesadaran, kompetensi, dan komitmen manusia
di dalamnya melalui arsitektur kesiapan yang terukur.

TISE adalah janji masa depan: sebuah paradigma rekayasa yang memastikan
bahwa di tengah kemajuan teknologi, tidak akan ada satu pun talenta
manusia yang terbuang sia-sia.

\section{Daftar Pustaka}\label{daftar-pustaka}

Morelli, N. (2002). Designing product/service systems: A methodological
exploration. \emph{Design Issues}, 18(3), 3--17.

Vezzoli, C., et al.~(2014). \emph{Product-Service System Design for
Sustainability}. Sheffield, U.K.: Greenleaf Publishing.

Berkovich, M., Leimeister, J. M., Hoffmann, A., \& Krcmar, H. (2014). A
requirements data model for product service systems. \emph{Requirements
Engineering}, 19, 161--186.

Sholihah, M., Mitake, Y., Nakada, T., \& Shimomura, Y. (2019).
Innovative design method for a valuable product-service system:
Concretizing multi-stakeholder requirements. \emph{Journal of Advanced
Mechanical Design, Systems and Manufacturing}, 13(5).

Giordano, F., Morelli, N., De Götzen, A., \& Hunziker, J. (2018). The
stakeholder map: A conversation tool for designing people-led public
services. \emph{ServDes2018 Conference Proceedings}.

Lievens, A., \& Blazevic, V. (2021). A service design perspective on the
stakeholder engagement journey during B2B innovation. \emph{Industrial
Marketing Management}, 95, 128--141.

\bookmarksetup{startatroot}

\chapter{Inovasi Pendidikan Teknik VALORAIZE melalui Paradigma
TISE}\label{inovasi-pendidikan-teknik-valoraize-melalui-paradigma-tise}

\section{Fondasi Filosofis: Paradigma Triune-Intelligence Smart
Engineering
(TISE)}\label{fondasi-filosofis-paradigma-triune-intelligence-smart-engineering-tise}

Pendidikan teknik masa depan menuntut lebih dari sekadar penguasaan
rumus; ia memerlukan kerangka kerja yang mampu menyinergikan kapasitas
teknologi dengan nilai kemanusiaan dan batasan alam. Paradigma
\textbf{Triune-Intelligence Smart Engineering (TISE)} mendefinisikan
rekayasa sebagai aplikasi kreatif pengetahuan ilmiah untuk menghasilkan
artefak (``Rancang-Bangun'') yang memanfaatkan kekuatan alam secara aman
dan terkendali guna memecahkan masalah penting manusia.

Filosofi TISE menggeser fokus dari optimasi teknis murni (seperti biaya
dan efisiensi) menuju penciptaan nilai holistik. Hal ini dicapai melalui
sinergi tiga pilar kecerdasan:

\textbf{Kecerdasan Manusia (Homocordium):} Bertindak sebagai kompas
moral, etika, dan penentu ``apa yang penting'' bagi kesejahteraan
manusia.

\textbf{Kecerdasan Buatan (Homodeus/Homologos):} Memberikan kekuatan
komputasi, logika, dan kemampuan belajar adaptif sebagai mitra berpikir.

\textbf{Kecerdasan Alam (Natural Intelligence):} Menetapkan batasan
fisik, hukum termodinamika, dan parameter keberlanjutan ekosistem.

Setiap artefak atau sistem yang lahir dari paradigma ini wajib memenuhi
enam karakteristik utama:

{\def\LTcaptype{none} % do not increment counter
\begin{longtable}[]{@{}
  >{\raggedright\arraybackslash}p{(\linewidth - 2\tabcolsep) * \real{0.5000}}
  >{\raggedright\arraybackslash}p{(\linewidth - 2\tabcolsep) * \real{0.5000}}@{}}
\toprule\noalign{}
\begin{minipage}[b]{\linewidth}\raggedright
Karakteristik
\end{minipage} & \begin{minipage}[b]{\linewidth}\raggedright
Deskripsi Teknis
\end{minipage} \\
\midrule\noalign{}
\endhead
\bottomrule\noalign{}
\endlastfoot
Strong & Memiliki Core Engine fundamental untuk konversi energi sumber
menjadi energi kerja. \\
Smart & Memiliki PUDAL Engine sebagai inti kognitif untuk persepsi dan
adaptasi berkelanjutan. \\
Extended Range & Mampu mengoptimalkan nilai di lima dimensi energi:
Product, Service, Knowledge, Value, dan Environmental. \\
Realistic & Kinerja artefak divalidasi secara empiris menggunakan
metodologi PICOC yang ketat. \\
Doable & Kompleksitas desain dikelola melalui dekomposisi hierarkis
struktur ASTF. \\
Methodic & Pengembangan mengikuti proses sistematis yang menghubungkan
desain dan validasi (V-Model). \\
\end{longtable}
}

-\/-------------------------------------------------------------------------------

\section{Arsitektur Platform VALORAIZE Learning: Dekomposisi
ASTF}\label{arsitektur-platform-valoraize-learning-dekomposisi-astf}

VALORAIZE mengadopsi kerangka kerja \textbf{ASTF} (\emph{Application,
System, Technology, Fundamental}) untuk membedah struktur pendidikan
teknik menjadi lapisan-lapisan yang dapat dikelola secara presisi.

\textbf{Layer Aplikasi (A):} Fokus pada masalah utama mahasiswa:
pencapaian N buah Capaian Pembelajaran Mata Kuliah (CPMK) dalam durasi
15 minggu. Solusinya berupa lingkungan belajar digital yang mewajibkan
penyerahan ``Karya'' sebagai bukti tertulis penguasaan materi.

\textbf{Layer Sistem (S):} Mengatur integrasi antara \emph{Knowledge
Marketplace}, Sistem Peta Pengetahuan Dasar (tempat berlatih kompetensi
mikro), dan Sistem Peta Aplikasi (tempat menyusun solusi kompleks).

\textbf{Layer Teknologi (T):} Menggunakan mesin kognitif berbasis AI
untuk mengelola transaksi pengetahuan dan platform digital untuk
mencatat akumulasi kompetensi mahasiswa.

\textbf{Layer Riset Fundamental (F):} Berpijak pada prinsip optimasi
jadwal waktu (riset operasi) dan teori konversi nilai PSKVE dalam proses
transformasi belajar menjadi aset intelektual.

\textbf{Aliran Vertikal:} Kerangka ASTF ini tidak berdiri sendiri; ia
bekerja selaras dengan V-Model. Desain diturunkan dari Layer A menuju
Layer F (sisi kiri V), sementara validasi dilakukan dari Layer F kembali
naik ke Layer A (sisi kanan V) untuk memastikan setiap prinsip
fundamental mendukung aplikasi akhir.

-\/-------------------------------------------------------------------------------

\section{Mekanisme Operasional: Siklus Kognitif PUDAL dan Transaksi
PSKVE}\label{mekanisme-operasional-siklus-kognitif-pudal-dan-transaksi-pskve}

VALORAIZE beroperasi sebagai ``mesin rekayasa pembelajaran'' melalui dua
siklus utama: kognisi dan konversi nilai.

\subsection{Siklus PUDAL (Proses
Kognitif)}\label{siklus-pudal-proses-kognitif}

Mahasiswa memproses setiap unit pengetahuan melalui lima tahap:

\textbf{Perceive:} Merasakan kebutuhan pengetahuan berdasarkan
\emph{Learning Outcome}.

\textbf{Understand:} Mencerna materi melalui pola dan hubungan antar
konsep.

\textbf{Decision-making:} Merancang rute belajar untuk memecahkan
persoalan CPMK.

\textbf{Act:} Mengeksekusi rute dengan menghasilkan ``Produk
Pengetahuan''.

\textbf{Learning:} Mengevaluasi umpan balik untuk memperbarui skema
kognitif.

\subsection{Siklus Konversi Energi PSKVE (Ekonomi
Pengetahuan)}\label{siklus-konversi-energi-pskve-ekonomi-pengetahuan}

VALORAIZE memandang proses belajar sebagai siklus konversi transaksional
lima dimensi energi:

\textbf{Knowledge (Energi Pengetahuan):} Algoritma, skema, dan data yang
dikuasai.

\textbf{Product (Energi Produk):} Karya nyata atau tugas yang
dihasilkan.

\textbf{Service (Energi Layanan):} Kapasitas memberikan solusi yang
memuaskan kebutuhan CPMK.

\textbf{Value (Energi Nilai):} ``Mata uang'' berupa nilai atau pengakuan
kompetensi.

\textbf{Environmental (Energi Lingkungan):} Keberlanjutan rute belajar
dan dampaknya pada ekosistem kelas.

Dalam siklus ini, \textbf{Value Energy} (nilai/kredit) yang diberikan
dosen digunakan untuk ``membeli'' atau menggerakkan \textbf{Knowledge
Energy} mahasiswa, yang kemudian dikonversi menjadi \textbf{Product
Energy} (tugas). Keberhasilan konversi ini mengisi ``Rekening Kekayaan
Mahasiswa,'' di mana nilai akhir adalah akumulasi dari kualitas
transaksi tersebut.

-\/-------------------------------------------------------------------------------

\section{Knowledge Marketplace dan Navigasi Rute
Belajar}\label{knowledge-marketplace-dan-navigasi-rute-belajar}

VALORAIZE memfasilitasi transaksi ``Produk Pengetahuan'' dalam sebuah
pasar digital yang transparan.

\subsection{Mekanisme Marketplace}\label{mekanisme-marketplace}

\textbf{Dosen:} Menayangkan ``Iklan Produk Pengetahuan'' yang
mendefinisikan \textbf{APA} (\emph{WHAT}) yang harus dicapai sebagai
profil target.

\textbf{Mahasiswa:} Mensubmit ``Produk Pengetahuan'' yang melaporkan
\textbf{BAGAIMANA} (\emph{HOW}) rute kompetensi tersebut dijalani.

\subsection{Navigasi 12-Langkah: Analogi Bandung ke
Stanford}\label{navigasi-12-langkah-analogi-bandung-ke-stanford}

Untuk memecahkan masalah kompleks, mahasiswa menyusun rute aplikasi
seperti perjalanan dari Bandung ke Stanford University:

Jalan Kaki (Rumah -\textgreater{} Gerbang). 2. Taksi (Gerbang
-\textgreater{} Stasiun Bandung). 3. Jalan Kaki (Platform
-\textgreater{} Gerbong). 4. Kereta Argo Parahyangan (Bandung
-\textgreater{} Gambir). 5. Jalan Kaki (Gerbong -\textgreater{} Platform
Gambir). 6. Jalan Kaki (Platform -\textgreater{} Bus Damri). 7. Bus
Damri (Gambir -\textgreater{} Terminal 3 Soetta). 8. Prosedur Bandara
(Check-in, Imigrasi, Boarding). 9. Terbang (Soetta -\textgreater{} SFO
via Taipei). 10. Prosedur SFO (Imigrasi, Bagasi). 11. Bus SamTrans (SFO
-\textgreater{} Palo Alto). 12. Stanford Shuttle (Palo Alto
-\textgreater{} Kampus).

Dalam setiap langkah, mahasiswa harus memilih moda:

\textbf{Jalan Kaki:} Melakukan penurunan matematis dan logis secara
manual (derivasi fundamental).

\textbf{Berkendaraan:} Menggunakan rumus, tabel, atau program komputer
(aplikasi praktis).

\subsection{Struktur Mandatori Produk Pengetahuan (20
Item)}\label{struktur-mandatori-produk-pengetahuan-20-item}

Setiap karya mahasiswa wajib mengikuti struktur standar untuk menjamin
kualitas:

Judul; 2. Subjudul; 3. Abstrak; 4. Pendahuluan; 5. Masalah; 6. Konsep
Solusi \& Kompetensi; 7. Laporan Kompetensi yang Didemonstrasikan; 8.
Studi Literatur; 9. Peta Dasar (Rute \& Kendaraan); 10. Peta Aplikasi;
11. Metodologi; 12. \emph{Gap} yang Dituju; 13. Analisis Tarik Mundur
(\emph{Backwards Analysis}); 14. Evaluasi Pilihan Rute; 15. Susunan Rute
Terpilih; 16. Hasil \& Diskusi; 17. Laporan Titik Antara; 18. Kesimpulan
\& Saran; 19. Lampiran; 20. Rujukan.

-\/-------------------------------------------------------------------------------

\section{Metodologi Validasi: V-Model dan PICOC
Berlapis}\label{metodologi-validasi-v-model-dan-picoc-berlapis}

VALORAIZE memastikan keterandalan pendidikan melalui validasi berbasis
bukti pada setiap lapisan ASTF di sisi kanan V-Model.

\textbf{V-Model:} Sisi kiri mendekomposisi kebutuhan (dari Aplikasi
hingga Riset Fundamental). Sisi kanan melakukan validasi dari bawah ke
atas.

\textbf{PICOC Berlapis:} Validasi dilakukan secara sistematis pada
setiap tingkatan:

\textbf{PICOC(F):} Memvalidasi teori konversi nilai (Misal: Apakah model
PSKVE benar-benar mengoptimalkan alokasi waktu?).

\textbf{PICOC(T):} Memvalidasi keandalan teknologi (Misal: Akurasi AI
dalam mendeteksi profil kompetensi).

\textbf{PICOC(S):} Memvalidasi integritas sistem (Misal:
\emph{Throughput} pencapaian CPMK dalam 15 minggu).

\textbf{PICOC(A):} Memvalidasi dampak pada pengguna (Misal: Kepuasan
mahasiswa dan dosen terhadap lingkungan belajar).

-\/-------------------------------------------------------------------------------

\section{Studi Kasus: Mata Kuliah Sinyal dan Sistem
(EL2007)}\label{studi-kasus-mata-kuliah-sinyal-dan-sistem-el2007}

Implementasi pada EL2007 mengubah silabus statis menjadi rute dinamis 15
minggu.

\textbf{Contoh Destinasi Antara (Minggu 2: Sistem LTI):}

\textbf{Target:} Mahasiswa menguasai Integral Konvolusi.

\textbf{Metodologi ``Walking'':} Mahasiswa wajib melakukan penurunan
matematis integral konvolusi secara manual dari prinsip impuls dasar
(Materi Oppenheim Bab 2).

\textbf{Metodologi ``Driving'':} Mahasiswa menggunakan fungsi
\texttt{scipy.signal.convolve} di Python untuk mensimulasikan respons
sistem terhadap sinyal masukan kompleks.

\textbf{Output:} Mahasiswa menyusun ``Produk Pengetahuan'' dengan 20
item mandatori, menyertakan \emph{Journal Belajar} yang menceritakan
proses berpikir dari titik buta (tidak paham konvolusi) hingga mampu
mendesain sistem filter.

Buku teks (Oppenheim \& Schaum's) diposisikan sebagai ``Peta Dasar'' dan
``Kendaraan'' yang harus dikendarai untuk mencapai target CPMK di minggu
ke-15.

\section{Perangkat Pendukung Teknis: Ontologi, Prolog, dan
Python}\label{perangkat-pendukung-teknis-ontologi-prolog-dan-python}

Keamanan dan akurasi platform VALORAIZE didukung oleh tiga pilar
teknologi:

\textbf{Ontologi:} Bertindak sebagai \textbf{Single Source of Truth
(SSOT)}. Seluruh definisi CPMK, hubungan antar kompetensi, dan syarat
kelulusan didefinisikan secara formal untuk menghindari ambiguitas
semantik.

\textbf{Prolog:} Digunakan untuk verifikasi formal kebijakan keamanan.
Mesin logika Prolog memastikan bahwa akses nilai dan perubahan data
``Rekening Kekayaan Mahasiswa'' mengikuti aturan yang tidak bisa
ditembus oleh konflik logika.

\textbf{Python:} Bahasa utama untuk simulasi rute belajar dan
pembangunan prototipe virtual. Python memungkinkan mahasiswa menguji
hipotesis teknis mereka dalam lingkungan yang aman sebelum dievaluasi
oleh dosen.

\section{Kesimpulan: Membentuk Insinyur sebagai
``Vokator''}\label{kesimpulan-membentuk-insinyur-sebagai-vokator}

Melalui paradigma TISE, VALORAIZE tidak sekadar melahirkan teknokrat,
melainkan seorang \textbf{Vokator}. Seorang Vokator adalah ``Arsitek
Masa Depan'' yang menyuarakan kebenaran melalui riset bertanggung jawab
dan produk pengetahuan yang bernilai.

Dengan menyatukan \emph{Homocordium} (etika), \emph{Homodeus} (logika
AI), dan \emph{Natural Intelligence} (alam), mahasiswa VALORAIZE dididik
untuk menjadi insinyur yang bijaksana, metodis, dan mampu menciptakan
``teater kehidupan yang megah'' bagi kemanusiaan. Inilah masa depan
pendidikan teknik yang holistik dan berkelanjutan.

\bookmarksetup{startatroot}

\chapter*{References}\label{references}
\addcontentsline{toc}{chapter}{References}

\markboth{References}{References}

\bookmarksetup{startatroot}

\chapter*{Lampiran}\label{lampiran}
\addcontentsline{toc}{chapter}{Lampiran}

\markboth{Lampiran}{Lampiran}

{\def\LTcaptype{none} % do not increment counter
\begin{longtable}[]{@{}
  >{\raggedright\arraybackslash}p{(\linewidth - 8\tabcolsep) * \real{0.1579}}
  >{\raggedright\arraybackslash}p{(\linewidth - 8\tabcolsep) * \real{0.1579}}
  >{\raggedright\arraybackslash}p{(\linewidth - 8\tabcolsep) * \real{0.2632}}
  >{\raggedright\arraybackslash}p{(\linewidth - 8\tabcolsep) * \real{0.1579}}
  >{\raggedright\arraybackslash}p{(\linewidth - 8\tabcolsep) * \real{0.2632}}@{}}
\toprule\noalign{}
\begin{minipage}[b]{\linewidth}\raggedright
Level
\end{minipage} & \begin{minipage}[b]{\linewidth}\raggedright
Domain
\end{minipage} & \begin{minipage}[b]{\linewidth}\raggedright
Core Engine
\end{minipage} & \begin{minipage}[b]{\linewidth}\raggedright
Function
\end{minipage} & \begin{minipage}[b]{\linewidth}\raggedright
Artifact
\end{minipage} \\
\midrule\noalign{}
\endhead
\bottomrule\noalign{}
\endlastfoot
Level 0 & Rekayasa Tradisional berbasis EMI (Energi, Materi, Informasi).
& Menyediakan tumpuan material, teknis, dan fisik serta fondasi
teknologi dan keandalan. & Core Engine (Intake, Encoder, Decoder,
Exhaust, Flywheel). & Konversi energi sumber menjadi kerja mekanis atau
teknis untuk menangani beban EMI secara statik atau dinamik. \\
Level 1 & Smart Engineering melalui konversi Energon. & Membangun
kapasitas kerja multidisiplin dan lingkungan kerja cerdas untuk
menghadapi beban masalah. & Smart Engine / Mesin PSKVE. & Mengonversi 5
jenis Energon (fisik, waktu, skill, token, peran) menjadi kapasitas
kerja terstruktur. \\
Level 2 & Rekayasa Agentik dan Triune Intelligence. & Memberikan
kecerdasan agentik, adaptasi, dan respon kontekstual pada artefak. &
Mesin PUDAL (Perception, Understanding, Decision, Action, Learning). &
Orkestrasi antara Natural, Collective, dan Artificial Intelligence
melalui Natural Language Prompting. \\
Level 3 & Identitas Naratif dan Teater Solusi. & Memberi makna,
motivasi, dan ruang bagi manusia untuk mementaskan misi hidupnya. &
Identitas Naratif dan Prompt Reflektif. & Menggali kisah hidup dan nilai
inti untuk diubah menjadi peran nyata dalam lingkungan yang
dirancang. \\
Level 4 & Meta-TISE dan DNA Desain. & Menyediakan pola pewarisan,
reproduksi desain, dan platform regeneratif. & Meta-TISE (DNA Desain dan
Platform TISE). & Pewarisan, adaptasi, dan ``perkawinan'' desain antara
DNA platform lama dengan gagasan baru. \\
Level 5 & Rekayasa Kemanusiaan dan Ekonomi Misi (VALORAIZE). &
Mobilisasi konfigurasi stakeholder bermisi untuk memecahkan persoalan
kemanusiaan. & Matrix of Mission dan MOS-7 (Mission Operating System). &
Pemetaan kontribusi timbal balik antar stakeholder (USER, SOURCE,
REGULATOR, PROVIDER) dalam siklus transformatif. \\
Level 6 & Ekosistem Keunggulan (Splendid Theaters of Life). & Panggung
global tempat seluruh manusia mementaskan talenta uniknya secara
maksimal. & Stadion Global Digital. & Amplifikasi talenta unik individu
menjadi dampak kemanusiaan dalam skala global. \\
\end{longtable}
}

\begin{center}\rule{0.5\linewidth}{0.5pt}\end{center}

\protect\phantomsection\label{refs}




\end{document}
